% [original page: 91-א]
ולכן:
\[ \frac{x_n\sqrt{a} + y_n\sqrt{b}}{2} = \left(\frac{x_1\sqrt{a} + y_1\sqrt{b}}{2}\right)^{2n-1} \]

אם נקח עתה בחשבון גם את כל הקומבינציות של הסימנים המופיעים ליד הערכים
$x$, $y$ במשוואה (1), הרי נקבל באופן סופי את הנוסחה

\[ \frac{x_n\sqrt{a} + y_n\sqrt{b}}{2} = \pm\left(\frac{x_1\sqrt{a} \pm y_1\sqrt{b}}{2}\right)^{2n-1} \]

בשביל כל $n$ טבעי: $n \geq 1$.

נוסחה זאת היא אנלוגית לנוסחה הידועה לגבי משוואת פל עצמה, ויש גם להבינה
כהשלמה ברוח הפירוש שניתן לגבי משוואת פל עצמה. % UNCERTAIN: exact wording, scan degraded

% [original page: 92]
\section*{פרק ז'.}

\subsection*{\S1. ההכללה הרביעית.}

\noindent\textbf{1. מבוא.}

בסעיפים הקודמים קבענו את יתר ההכללות למשוואת פל, ביחס למקדמי התבנית
ולמייצגים. % UNCERTAIN wording in this sentence

נרחיב עתה את הדברים עד כמה שאפשר בכדי להכליל את משוואת פל גם ביחס
למעריכים. אמנם על-ידי כך אנו עוברים כבר לטיפול במשוואות מסדר גבוה, או
בתבניות שאינן ריבועיות, נושא החורג ממסגרת חלק זה; אבל כפי שנראה יתכן
סוג מיוחד של משוואות מסדר גבוה אשר כל התורה העיקרית אליהן ניתנת לביסוס
באמצעים אלמנטריים אך ורק על סמך המסקנות שקיבלנו. % UNCERTAIN wording, dense typewriter text

מאידך כדאי לציין, שכל התיזה בהיקפה המלא, על דבר משוואות מהסוג הראוי
לכנותן ביחס למשוואת פל, שייכת לתורת המספרים האלגבראיים, ובייחוד לענף
המעניין -- לתורת היחידות של דיריכלה; אלא שכפי שכבר ציינו חורגים המקרים
המיוחדים בהם אנו מטפלים כאן מהמסגרת הכללית שנוצרה על-ידי דיריכלה,
והמסקנות מתפצלות בהרבה בהרבה בכל אחד מהמקרים דנן. % UNCERTAIN: several words illegible, best-effort reading

המשוואות יבחנו להבא הן מהצורות:

\[ ax^{2\mu} - by^{2\nu} = \pm 1 , \pm 2 , \pm 4 \]

במקום $\mu, \nu$ הם מספרים טבעיים כלשהם. % UNCERTAIN

את תכונות המשוואות הללו ניסו לחקור אנשים רבים ובעניינם נמצאים אחדים
מהמפרסמים של דורנו האחרונים. גם מהם יוחלץ חלקנו בידם, אבל יש גם ביניהם
כאלה שהצליחו לגלות הרבה דברים מעניינים, חומר על כך נביא בחומר
הביבליוגרפי המיוחד. % UNCERTAIN wording

שתי המסקנות החשובות ביותר שנתגלו המובאות בפרק זה הן:

א) קביעת שיטה לכתרון המשוואות דלעיל;

% [original page: 93]
ב) העובדה החשובה והמעניינת שכל פתרון $(x,y)$ של המשוואה מהצורות
דלעיל ניתן להתבטא בעזרת הפתרון המינימלי המינימלי של אותה משוואה
בצורה:

\[ x\sqrt{a} \pm y\sqrt{b} = \pm(x_1\sqrt{a} \pm y_1\sqrt{b})^N \]

\underline{בשביל $N$ מספרים.}

מלבד יתר המסקנות הללו נקבעה עוד עובדה שלישית לא פחות מעניינת בדבר
הקשר הקיים בין הפתרונות של סוג המשוואות מעבר, שגם הן יוצרות
\underline{כל, מהצורות:}
\[ ax^2 - by^{2\nu} = \pm 1 , \pm 2 , \pm 4 \]
\[ ax^{2\mu} - by^2 = \pm 1 , \pm 2 , \pm 4 \]

\noindent\textbf{2.} נטפל עתה במשוואה מהצורה

\[ (1) \qquad ax^{2\mu} - by^{2\nu} = \pm 1 \]

במקום ש-$\mu, \nu$ הם מספרים טבעיים שונים או שווים כלשהם, מאליו יוצא,
שיתרון המשוואה הנ"ל במספרים $x, y$ כלמים הוא שקול כנגד פתרון שתי
המשוואות:

\[ (2) \qquad ar^2 - bs^2 = \pm 1 , \qquad r = x^{\mu} , \quad s = y^{\nu} \]

\underline{בבת-אחת.}

כי קיום המערכת האחרונה במספרים $r,s,x,y$ שלמים גורר אחריו את פתרון
המשוואה (1) והדברים ניתנים גם להיפוך.

לפיכך נוכל לשטל בכל בחקירותינו מיד במערכת (2) במקום המשוואה (1).

\noindent\textbf{3.} נניח עתה עוד, שהמשוואה (1) איננה מהצורה המצומצמת המיוחדת:

\[ (3) \qquad r^2 - Ds^2 = 1 , \qquad r = x^{\mu} , \qquad s = y^{\nu} \]

למקרה המיוחד הזה נחזור עוד לטפל בו בסעיפים המיוחדים עליו.

% [original page: 94]
מן המסקנות שנתקבלו לגבי ההכללות של משוואת פל הראשונות אנו יודעים
כבר, כי בנוסחה

\[ (3') \qquad r_n\sqrt{a} \pm s_n\sqrt{b} = \pm(r_1\sqrt{a} \pm s_1\sqrt{b})^{2n-1} \]
\[ n = 1, 2, 3, \ldots \]

במקום ש-$r_1, s_1$ הנן הפתרון המינימלי של המשוואה הנ"ל נכללים כל
פתרונות המשוואה

\[ (4) \qquad ar^2 - bs^2 = 1 . \]

להבא נזניח את כל הקומבינציות האפשריות לגבי חילופי הסימנים המופיעים
בנוסחה, כי אינן משפיעות על תוצאות החקירות.

מטרתנו היא, אפוא, עתה לקבוע $n$ מתאים שבשבילו יתקיימו התנאים:

\[ (5) \qquad r_n = x^{\mu} , \qquad s_n = y^{\nu} \]

\underline{בבת-אחת}; במלים אחרות: עלינו למצוא פתרון למשוואה (4) שבו
$r, s$ הם חזקות מהסדרים $\mu, \nu$ בהתאמה.

\noindent\textbf{4.} נסמן לשם הקיצור

\[ (6) \qquad 2n - 1 = m \]

ונפתח את האגף הימני שבנוסחה (3') לפי נוסחת ההיתוח של בינום ניוטון:
לאחר הפרדת אי-הרציונליות הזוגות לחוד והאחרות מכל אגפי החיוורון נקבל

\[ r_n = r_1^{m}(\sqrt{a})^m + B_2^{(m)} r_1^{m-2} s_1^2 b(\sqrt{a})^{m-2} + \cdots
+ B_{m-1}^{(m)} r_1 s_1^{m-1} b^{\frac{m-1}{2}}\sqrt{a} \; ; \]

\[ s_n = B_1^{(m)} r_1^{m-1} s_1(\sqrt{a})^{m-1} b + B_3^{(m)} r_1^{m-3} s_1^3(\sqrt{a})^{m-3}(\sqrt{b})^3
+ \cdots + B_m^{(m)} s_1^m(\sqrt{b})^m \]

במקום ש-$B_k^{(m)}$ הם כל המקדמים הבינומיאליים בגבול

\[ 1 \leq k \leq m \]

% [original page: 95]
\noindent מכאן אנו מקבלים:

\[ (7) \quad
\begin{cases}
r_n = r_1\left(aPr_1 + m s_1^{m-1} b^{\frac{m-1}{2}}\right) \\[4pt]
s_n = s_1\left(bQs_1 + m r_1^{m-1} a^{\frac{m-1}{2}}\right)
\end{cases}
\]

במקום ש-$P, Q$ הם מספרים שלמים מסויימים, בשים לב לכך, כי:

\[ (8) \qquad B_{m-1}^{(m)} = B_1^{(m)} = m \]

\noindent\textbf{5.} נניח עתה שהערכים המינימליים $r_1, s_1$ אינם חזקתיים המקיימים
את המערכת (5). מוכרח, אפוא, להיות, לחזור:

\[ (9) \quad
\begin{cases}
r_1 = r'^{\mu} r''^{\mu - \lambda} \\
s_1 = s'^{\nu} s''^{\nu - \delta}
\end{cases}
\]

במקום ש-$r', r'', s', s'', \lambda, \delta$ הם מספרים שלמים מסויימים;
כמובן יכול גם להיות לפירוח: % ILLEGIBLE: word uncertain

\[ (10) \qquad \lambda \leq \mu , \qquad \delta \leq \nu \]

ו-:

\[ (11) \qquad r' \geq 1 , \qquad s' \geq 1 \]

במקרה של יוריון % ILLEGIBLE: "יוריון"/"שיוויון" unclear
במערכת (10) הפתרון המינימלי $r_1, s_1$ מקיים את התנאים (5). במקרה
של יוריון במערכת (11) אין כל גורמים חזקתיים מסדר $\mu, \nu$ לערכים
$r_1, s_1$ בהתאמה.

\noindent\textbf{6.} נפרק עתה את $r'', s''$ לגורמים ראשוניים; ננית שההצגות המתקבלות הן:

\[ (12) \quad
\begin{cases}
r'' = p_1^{\alpha_1} p_2^{\alpha_2} p_3^{\alpha_3} \cdots p_k^{\alpha_k} \\
s'' = q_1^{\beta_1} q_2^{\beta_2} q_3^{\beta_3} \cdots q_1^{\beta_1}
\end{cases}
\]

במקום שהמספרים

\[ p_1, p_2, p_3, \ldots, p_k \; ; \; q_1, q_2, q_3, \ldots, q_1 \]

% [original page: 96]
\noindent הם ראשוניים ושונים בינהם.

\underline{למד"ז יהיה:}

\[ (13) \quad
\begin{cases}
r''^{\mu-\lambda} = p_1^{\zeta_1} p_2^{\zeta_2} p_3^{\zeta_3} \cdots p_k^{\zeta_k} \\
s''^{\nu-\delta} = q_1^{\xi_1} q_2^{\xi_2} q_3^{\xi_3} \cdots q_1^{\xi_1}
\end{cases}
\]

\underline{במקום}

\[ (14) \quad
\begin{cases}
\zeta_1 = (\mu-\lambda)\alpha_1 , \quad \zeta_2 = (\mu-\lambda)\alpha_2 , \; \ldots \\
\zeta_k = (\mu-\lambda)\alpha_k \; ; \\[4pt]
\xi_1 = (\nu-\delta)\beta_1 , \quad \xi_2 = (\nu-\delta)\beta_2 , \; \ldots \\
\xi_1 = (\nu-\delta)\beta_1 .
\end{cases}
\]

מהמשוואה (1) נובע מיד, כי:

\[ (15) \qquad (ar_1 , bs_1) = 1 \]

ולכן נסיק גם מיד מהמערכת (7), שאם $r_n, s_n$ מקיימים את התנאים (5)
מוכרח להיות:

\[ (16) \qquad m = M p_1^{\vartheta_1} p_2^{\vartheta_2} p_3^{\vartheta_3} \cdots p_k^{\vartheta_k}
= N q_1^{\tau_1} q_2^{\tau_2} q_3^{\tau_3} \cdots q_1^{\tau_1} \]

בשביל מספרים שלמים $M, N$ מסויימים, במקום שמתקיימים:

\[ (17) \quad
\begin{cases}
\vartheta_1 + \zeta_1 \equiv 0 , \quad \vartheta_2 + \zeta_2 \equiv 0 , \; \ldots \; , \\
\vartheta_k + \zeta_k \equiv 0 \pmod{\mu}
\end{cases}
\]

מחד גיסא, וגם:

\[ (18) \quad
\begin{cases}
\tau_1 + \xi_1 \equiv 0 , \quad \tau_2 + \xi_2 \equiv 0 , \; \ldots \; , \\
\tau_1 + \xi_1 \equiv 0 \pmod{\nu}
\end{cases}
\]

מאידך גיסא.

אולם בשים לב למערכת הקודמים (14) נוכל לשטל את התנאים האחרונים עד

% [original page: 97]
\noindent שנקבל תנאים יותר נוחים, והם:

\[ (19) \quad
\begin{cases}
\vartheta_1 - \lambda\alpha_1 \equiv 0 , \quad \vartheta_2 - \lambda\alpha_2 \equiv 0 , \; \ldots \\
\vartheta_k - \lambda\alpha_k \equiv 0 \pmod{\mu}
\end{cases}
\]

מחד גיסא, ו-

\[ (20) \quad
\begin{cases}
\tau_1 - \delta\beta_1 \equiv 0 , \quad \tau_2 - \delta\beta_2 \equiv 0 , \; \ldots \\
\tau_1 - \delta\beta_1 \equiv 0 \pmod{\nu}
\end{cases}
\]

מאידך גיסא.

הלאה, מתני החקרים (16) לגבי הערך $m$ שבשבילו מתקיימת המערכת (5)
אנו מקבלים את הצורה:

\[ (21) \qquad m = p_1^{\vartheta_1}p_2^{\vartheta_2}p_3^{\vartheta_3}\cdots p_k^{\vartheta_k}
\cdot q_1^{\tau_1}q_2^{\tau_2}q_3^{\tau_3}\cdots q_1^{\tau_1} \cdot M_1 \]

במקום ש-$M_1$ הנו מספר טבעי מסויים מכאן

\noindent\textbf{\underline{משפט 1.}}

אם בשביל המעריך $m = 2n-1$ מתקיימת המערכת (5) אזי מוכרח להיות
$m$ מהצורה (21).

\noindent\textbf{7.} המסקנה שקבלנו בסעיף הקודם מאפשרת לנו לקבוע שיטה, שלפיה אפשר
תמיד לחזור% ILLEGIBLE: perhaps "לבחור"
את המשוואה (1) בכל המקרים, אם כי שיטה זאת אינה דיה
כשלעצמה בכדי לקבוע מראש אם יש בכלל פתרון בכלל למשוואה (1).

ואמנם אם נסמן

\[ (22) \qquad N_1 = p_1^{\vartheta_1}p_2^{\vartheta_2}p_3^{\vartheta_3}\cdots p_k^{\vartheta_k}
\cdot q_1^{\tau_1}q_2^{\tau_2}q_3^{\tau_3}\cdots q_1^{\tau_1} \]

אזי נוכל למצוא את ערכו של $N_1$ הואיל וערכיהם של $r,s$ ידועים
מראש על-ידי תורת המשוואה (4). לפיכך נוכל לקבוע מיד, כי:

\[ (23) \qquad r_n\sqrt{a} + s_n\sqrt{b} = (r_1\sqrt{a} + s_1\sqrt{b})^{N_1 M_1}
= (r_{N_1}\sqrt{a} + s_{N_1}\sqrt{b})^{M_1} . \]

% [original page: 98]
\noindent אם יתברר עתה שהערכים $r_{N_1}, s_{N_1}$ עצמם אינם מקיימים עדיין את
התנאים (5) נוכל לחזור על כל התהליך הקודם ביחס ל-$r_{N_1}, s_{N_1}$ ונקבל

\[ (24) \qquad M_1 = N_2 M_2 \]

במקום שערכו של $N_2$ ידוע, וכך הלאה, עד שלבסוף נגיע לידי מערך

\[ (25) \qquad N_u = M_u \]

שבשבילו הערכים $r_{N_u}, s_{N_u}$ מקיימים את התנאים (5) בכל מקרה
שיש רק פתרון למשוואה (1).

אלא, שלעולם אין לדעת מראש אם השיטה הזאת תוביל סוף סוף לידי פתרון;
כי יתכן הרי שאין פתרון בכלל למשוואה (1) הנתונה, כלומר שהמקדמים
הם $a, b$ והמעריכים $\mu, \nu$.

\noindent\textbf{8.} אם כי השיטה עצמה לקביעה בזה קווריה % ILLEGIBLE: "קווריה" unclear
מאפשרת לקבוע מראש אם יש
בכלל פתרון למשוואה (1) הרי ימצא פתרון אחר בצדה. ממנה נובעת עובדה
חשובה עד למאוד והמעניינת בפני עצמה, והיא:

\noindent\textbf{\underline{משפט 2.}}

אם $(x_1, y_1)$ הנו הפתרון המינימלי של המשוואה (1) אזי כל פתרון
$(x, y)$ אחר של אותה משוואה ניתן להתבטא בצורה:

\[ (26) \qquad x\sqrt{a} \pm y\sqrt{b} = \pm(x_1\sqrt{a} \pm y_1\sqrt{b})^N \]

\underline{בשביל $N$ מספרים.}

ההוכחה נובעת מעצם התהליך שעליו בנויה השיטה דלעיל. כירון % ILLEGIBLE: "כירון" unclear, perhaps "כיוון"
שלפי שיטה
זאת נוכל להגיע לידי פתרון $r_n, s_n$ המקיימים את המערכת (5) אך ורק
על-ידי מעבר של סורם % ILLEGIBLE: "סורם" unclear, perhaps "סורג"/"טור"
מעריכים (21)(24)(25).

% [original page: 99]
\subsection*{\S2. מקרים מיוחדים.}

\noindent\textbf{1.} נעבור עתה לתקורת % ILLEGIBLE: "לתקורת" unclear, perhaps "לחקירת"
אפשרויות מהצורות המיוחדות:

\[ (27) \qquad ax^{2\mu} - by^2 = \pm 1 , \qquad ax^2 - by^{2\nu} = \pm 1 \]

בתנאים הידועים.

משוואה מצורה כזאת היא למעשה מקרה מיוחד של המשוואה מהצורה הכללית
שחקרנו כבר קודם, לאחר שנניח

\[ \nu = 1 \qquad ; \qquad \mu = 1 \]

אולם כפי שיתברר קיים מצב מיוחד ביחס לפתרונותיה של כל אחת מהמשוואות
האחרונות.

\noindent\textbf{2.} יהי $n_1$ אחד המעריכים אשר בשבילם מקיימים $r_m, s_m$ את המערכת
(5) בסעיף הקודם.

יהי עוד

\[ (28) \qquad n_1 = p_1^{\alpha_1}p_2^{\alpha_2}p_3^{\alpha_3}\cdots p_k^{\alpha_k} \]

הפירוק של $n_1$ לגורמים ראשוניים. אפשר אז לתאר את $n_1$ בצורה

\[ (29) \qquad n_1 = N_1 p_k^{\alpha_k} \]

במקום ש-$N_1$ הנו מספר טבעי מסויים. על סמך המסקנות והרעיונות
הידועים כבר נוכל לקבוע את הנוסחאות:

\[ (30) \quad
\begin{cases}
r_m = x^{\mu} = r_{N_1}\bar{r} = r_{N_1}\left(aP'r_{N_1} + p_k^{\alpha_k}s_{N_1}^{p_k^{\alpha_k}-1}
b^{\frac{p_k^{\alpha_k}-1}{2}}\right) \\[6pt]
s_m = y^{\nu} = s_{N_1}\bar{s} = s_{N_1}\left(bQ's_{N_1} + p_k^{\alpha_k}r_{N_1}^{p_k^{\alpha_k}-1}
a^{\frac{p_k^{\alpha_k}-1}{2}}\right)
\end{cases}
\]

בשביל $P', Q'$ שלמים מסויימים.

יהי עתה:

% [original page: 100]
\[ (31) \quad
\begin{cases}
r_{N_1} = r'^{\mu} r''^{\mu-\lambda} \\
s_{N_1} = s'^{\nu} s''^{\nu-\delta}
\end{cases}
\]

בכבל $\lambda, \delta$ מסויימים, כך ש-$r', r'', s', s''$ הם מספרים
טבעיים. הצגה כזאת מוכרחת להתקיים כפי שכבר הסברנו.

בדרך שהלכנו בה בסעיף הקודם נווכח מיד ש-$r_m = x^{\mu}, s_m = y^{\nu}$
יהוו חזקות של $r_{N_1}, s_{N_1}$ בהתאמה רק אז כאשר $p_k^{\alpha_k}$
יתחלק בגורמים הראשוניים של $r', s'$. אבל $p_k$ הנו מספר ראשוני
מצד אחד ומצד שני:

\[ (r_{N_1}, s_{N_1}) = 1 \; ; \]

לפיכך נסיק מיד, שאחד העריכים $r_{N_1}, s_{N_1}$ חייב להיות בעצמו
מוחזק מ-$\mu$ או $\nu$ בהתאמה. מכאן אנו מקבלים את

\noindent\textbf{\underline{משפט 3.}}

% [original page: 101]
בשביל כל $N_1$ המקיים את התנאים:

\[ N_i \, p_i^{m_j} = n_1 \]
\[ 1 \leq i \leq k \]
\[ 1 \leq m_j \leq \alpha_i . \]

% [original page: 102]
\subsection*{\S3. מקרים מיוחדים נוספים.}

\noindent\textbf{1.} בסעיף הקודם ציינו מקרים מיוחדים ביחס למעריכים. בסעיף זה נראה
שאפשר לקבל מסקנות מיוחדות במקרים שונים ביחס למקדמים. בתורת המשוואות
הדיופנטיות נודעת חשיבות רבה להוכחת אי-אפשרות הפתרון של משוואות
מסויימות. ידועה הבעיה המפורסמת של פרמה; אולם בעקבותיה של זאת
האחרונה התפתחו הרבה בעיות קרובות לה במידת מה או גם כאלו שנתקלים
בהן על דרך החיפושים אחרי פתרון בעיית פרמה. בפרק מיוחד נעמוד עוד
על הקשר בין המשוואות המכלליות של כל שחקרנו אותן בחלק זה לבין בעיית
פרמה בכללה. % UNCERTAIN wording, dense typewriter text

אולם כאן נביא רק כמה דוגמאות פשוטות למדי.

\noindent\textbf{\underline{משפט 1.}}

המשוואה הדיופנטית

\[ (1) \qquad x^{2n} - 2y^n = 1 , \qquad n > 1 \]

היא בלתי-אפשרית במספרים $x, y$ שלמים.

\noindent\textbf{\underline{הוכחה.}}

מהמשוואה (1) נובע מיד, כי:

\[ (2) \qquad (x^n - 1)(x^n + 1) = 2y^n \]

אולם

\[ (3) \qquad (x^n - 1 , x^n + 1) = 2 \]

סכיוון שממשוואה (1) רואים מיד כי $x$ הוא לא-זוגי, ולכן המספרים
$x^n - 1, x^n + 1$ הם שניהם זוגיים, ומלבד המספר 2 אין להם כל מחלק
משותף אחר, כפי שנובע מ-(1).

מ-(2) אפשר להסיק מיד על סמך הפירוק החד-ערכי, שמתקיימים בבת-אחת:

\[ (4) \qquad x^n + 1 = 2\zeta^n , \qquad x^n - 1 = \varepsilon^n \]

% [original page: 103]
במקום

\[ (5) \qquad \varepsilon\zeta = y \]

אולם מבלי להגביל את עצמנו אפשר להניח, כי:

\[ y > 0 \]

ולכן גם

\[ (6) \qquad \varepsilon > 0 , \qquad \zeta > 0 . \]

מהמשוואה החניונית של המערכת (4) מקבלים:

\[ \varepsilon^n - x^n = -1 \]

ומכאן:

\[ (7) \qquad (\varepsilon - x)P = -1 \]

בשביל מספר $P$ שלם המקיים:

\[ (8) \qquad P = \varepsilon^{n-1} + \varepsilon^{n-2}x + \varepsilon^{n-3}x^2 + \cdots + x^{n-1} \]

מהחיוורון (7) נובע מיד, כי:

\[ (9) \qquad P = \pm 1 \]

אבל מ-(8) אפשר להיווכח, כי:

\[ (10) \qquad P > 1 \; ; \]

יחדיו מבלי להגביל את עצמנו אפשר להניח, כי: $x > 0$; ולפי (6) גם
$\varepsilon > 0$ ולכן בשביל כל $n > 1$ יהיה:

\[ P \geq 1 + 1 = 2 \]

חני החיוורונים (9)(10) סותרים זה את זה ועל כן המשוואה (1) הייתה
בלתי אפשרית. % UNCERTAIN: first word of sentence unclear

\noindent\textbf{מ.ש.ל.}

על סמך משפט זה נוכל להוכיח מיד עוד חיובות % ILLEGIBLE: "חיובות" unclear
רבה ויחידה לבעיית פרמה, כפי שיתברר עוד במקום אחר.

נביא את המשפט כאן, אף כי למעשה היה צריך להקדים אותו קודם פרק מיוחד

% [original page: 104]
\noindent לחקירת עם המשוואה הבה דנים במשפט; אולם הוכחת העובדה היסודית המוצאת
את ביסודה במשפט זה מתחמכת על הדברים שכבר הובאו בחלק זה של תורת
התבניות הבינוניות. % UNCERTAIN wording

\noindent\textbf{\underline{משפט 2.}}

אם יש פתרון למשוואה

\[ (1) \qquad u^2 - Dv^2 = 4 , \qquad (u, v) = 1 \]

בשביל $D$ מסויים, וחוץ מזה יש גם פתרון למשוואה

\[ (2) \qquad s^2 - 3^{2(n-1)}Dt^2 = 4 , \qquad (s, t) = 1 \]

בשביל

\[ n > 1 \]

אזי אין פתרון למשוואה

\[ (3) \qquad x^{2n} - Dy^{2n} = 1 , \]
\[ (4) \qquad n > 1 \]

במספרים $x, y$ שלמים.

\noindent\textbf{\underline{הוכחה.}}

כידוע יש תמיד פתרון למשוואה כל מהצורה

\[ (5) \qquad \varepsilon^2 - D\eta^2 = 1 . \]

ולפי פוארקדה % ILLEGIBLE: reference name/term unclear
קיים קשר בין פתרונות שתי המשוואות (1)(5) מהצורה:

\[ (6) \qquad \varepsilon_1 + \eta_1\sqrt{D} = \left(\frac{u_1 + v_1\sqrt{D}}{2}\right)^3 \]

במקום ש-$\varepsilon_1, \eta_1 ; u_1, v_1$ הם הפתרונות המינימליים של שתי
המשוואות הנ"ל.

מהקשר האחרון מקבלים בדרך כלל

\[ (7) \qquad \varepsilon_k + \eta_k\sqrt{D} = \left(\frac{u_1 + v_1\sqrt{D}}{2}\right)^{3k} \]

% [original page: 105]
\noindent לגבי כל פתרון $\varepsilon_k, \eta_k$ של המשוואה (5); או:

\[ (8) \qquad \varepsilon_k + \eta_k\sqrt{D} = \left(\frac{u_k + v_k\sqrt{D}}{2}\right)^3 \]

במקום ש-$u_k, v_k$ הנן פתרון מסויים של המשוואה (1), כפי שנובע
מהנוסחאות הידועות לפתרונות משוואת פל.

מהקשר האחרון מקבלים מיד:

\[ (9) \qquad \varepsilon_k + \eta_k\sqrt{D} =
\frac{u_k(u_k^2 + 3v_k^2 D) + v_k(3u_k^2 + v_k^2 D)\sqrt{D}}{8} \]

אולם

\[ u_k^2 - Dv_k^2 = 4 \]

ולכן

\[
\begin{cases}
u_k^2 + 3Dv_k^2 = 4(Dv_k^2 + 1) \\
3u_k^2 + v_k^2 D = 4(u_k^2 - 1)
\end{cases}
\]

ולפיכך נקבל ע"י השוואת החלקים הרציונליים והאי-רציונליים בשווירון
(9) את הקשרים הבאים:

\[ (10) \quad
\begin{cases}
\varepsilon_k = \dfrac{u_k(Dv_k^2 + 1)}{2} \\[6pt]
\eta_k = \dfrac{v_k(u_k^2 - 1)}{2}
\end{cases}
\]

נניח עת שיש פתרון למשוואה (3). ברור אז שיתקיימו הקשרים הקודמים:

\[ (11) \qquad x^n = \varepsilon_k , \qquad y^n = \eta_k \]

בשביל $k$ מסויים. לפיכך מוכרח להיות:

\[ (12) \quad
\begin{cases}
x^n = \dfrac{u_k(Dv_k^2 + 1)}{2} \\[6pt]
y^n = \dfrac{v_k(u_k^2 - 1)}{2}
\end{cases}
\]

מאידך נובע מהמשוואה

\[ u_k^2 - Dv_k^2 = 4 , \]

% [original page: 106]
\noindent כי:

\[ Dv_k^2 + 1 = u_k^2 - 3 , \qquad u_k^2 - 1 = Dv_k^2 + 3 \]

ולכן נסיק מיד, כי:

\[ (u_k , Dv_k^2 + 1) = 3^{\gamma} , \qquad \gamma = 1, 0 \]
\[ (v_k , u_k^2 - 1) = 3^{\lambda} , \qquad \lambda = 1 - \gamma , \qquad 0 \]

מכיוון ש:

\[ (u_k , v_k) \neq 3 . \]

נניח עתה שאין פתרון למשוואה

\[ s^2 - 3^{2(n-1)}Dt^2 = 4 ; \]

לפי"ז אין גם פתרון למשוואה מהצורה

\[ 3^{2(n-1)}u'^2 - Dv'^2 = 4 \]

כפי שקל להוכיח. מכאן נוכרח % ILLEGIBLE: "נוכרח" likely "מוכרח"
להיווכח שמתקיים בדרך כלל

\[ 3 \nmid u_k , \qquad 3 \nmid v_k \]

ולפיכך מתקיימים:

\[ (13) \qquad (u_k , Dv_k^2 + 1) = 1 . \]

על סמך הפירוק החד-ערכי נסיק מיד מהקשרים (12), כי:

\[ (14) \qquad u_k = P^n , \qquad Dv_k^2 + 1 = 2Q^n \]

מכיוון ש-$u_k$ הוא לא-זוגי, וגם

\[ (15) \qquad v_k = R^n , \qquad u_k^2 - 1 = 2S^n \]

במקום ש-$P, Q, R, S$ הם מספרים שלמים מסויימים. אבל מכאן נובע מיד,
כי:

\[ P^{2n} - 2S^n = 1 \]

מה שאינו יתכן, לפי משפט 1. ובזה הוכח המשפט.

% [original page: 107]
\section*{ביבליוגרפיה}

\foreignlanguage{english}{C.F. Gauss, \textit{Disquisitiones arithmeticae}, Art. 1--31; 94--215.}

\foreignlanguage{english}{L. Dirichlet, \textit{Werke} II, I.}

\foreignlanguage{english}{Dedekind, \textit{Zahlentheorie}.}

\foreignlanguage{english}{Hardy \& Wright, \textit{The Theory of Numbers}, p. 67, 93.}

\foreignlanguage{english}{W. Sierpinski, \textit{Teoria liczb}.}

\foreignlanguage{english}{B. Nieweglowski, \textit{Wiad. Mat.}, Warsaw, 12, 1908, str. 1--26.}

\foreignlanguage{english}{H. Weber, \textit{Math. Annalen} Bd. 43 (1893), 184--195.}

\foreignlanguage{english}{Degen, \textit{Canon Pellianus}, Kopenhagen 1817.}

\foreignlanguage{english}{A. Meyer, \textit{Diss. Z\"urich}, 1871; \textit{Vierteljahrschrift Naturf. Gesell. Z\"urich}, 32, 1887, 363--382.}

\foreignlanguage{english}{H. Poincar\'e, \textit{Comptes Rendus}, Paris, 91, 1880, 846.}

\foreignlanguage{english}{J. Perott, \textit{J. f. M.}, 96, 1884, 335--7.}

\foreignlanguage{english}{G. Eisenstein, \textit{J. f. M.}, 27, 1884, 86.}

\foreignlanguage{english}{Roberts, \textit{Proc. London Math. Soc.}, 9, 1887--8, 194.}

\foreignlanguage{english}{R. Remak, \textit{J.f.M.} 143, S. 106.}

\foreignlanguage{english}{D. Hilbert, \textit{G\"ottinger Nachrichten (Math.)}, 1897, 52--54.}

\foreignlanguage{english}{K. Peter, \textit{Casopis} 56, (1927), S. 57--66.}

\foreignlanguage{english}{P. Seeling, \textit{Archiv Math. Phys.}, 52, 1871, 40--9.}

\foreignlanguage{english}{M.A. Stern, \textit{J.f.M.} 10, 1833, 1--22, 154--156, 241--274; 11, 1834, 33--66, 142--168, 277--306, 311--350.}

% [original page: 108]
\foreignlanguage{english}{F. N. C. Egen, \textit{Handbuch der allgemeinen Arith.}, Berlin, 1819--20; ed. 2, I, 1833, 457; II, 1834, 467; ed. 3, I, 1846, 456; II, 1849, 468.}

\foreignlanguage{english}{G. Speckmann, \textit{Archiv Math. Phys.}, (2), 13, 1895, 327--333; 14, 1895, 443--5.}

\foreignlanguage{english}{A. Palmstr\"om, \textit{Bergens Museums \r{A}rbog for 1896}, Bergen, 1897, 14, 11 pp.}

\foreignlanguage{english}{A. Cunningham \& R. W. D. Christie, \textit{Math. Quest. Educ. Times}, 73, 1900, 115--7.}

\foreignlanguage{english}{J. Ricalde, \textit{L'interm\'ediaire des math.}, 8, 1901, 256; 6, 1899, 51.}

\foreignlanguage{english}{G. Fontan\'e, \textit{Bull. math. \'el\'ementaire}, 15, 1909, 178--182.}

\foreignlanguage{english}{O. Perron, \textit{J. f. M.}, 1, 1914, 1--73.}

\foreignlanguage{english}{Legendre, \textit{M\'em. Acad. Sc. Paris}, 1785, 549--551; \textit{Th\'eorie des nombres}, 1798, 65--67; ed. 3, 1, 1830, 64--71.}

\foreignlanguage{english}{Jacobi, \textit{Monatsber. Akad. Wiss.}, Berlin, 1837, 127; \textit{J. f. M.} 30, 1845, 166; \textit{Werke} VI, 263--4; \textit{Opuscula Mathematica} 1, 1846, 324--5.}

\foreignlanguage{english}{E. de Jonqui\`eres, \textit{Comptes Rendus}, Paris, 126, 1898, 863--871, 991--7.}

\foreignlanguage{english}{Carmichael, R. D., \textit{Zentralblatt}, 16, 1937, S. 10.}

\foreignlanguage{english}{W. Ljunggren, \textit{Zentralblatt} 16, S. 8.}

% [original page: 109]
\foreignlanguage{english}{V. Tartakowski, \textit{Bull. Acad. Sc. de l'U.R.S.S.} 1926, S. 301--324.}

\foreignlanguage{english}{Axel Thue, \textit{Skrifter Videnskapsselskapet}, Kristiania 1909.}

\foreignlanguage{english}{T. Nagell, \textit{L'analyse ind\'etermin\'ee de degr\'e sup\'erieur}, M\'emorial des sc. math. fasc. XXXIX, S. 32--3, 54--5.}

\foreignlanguage{english}{S. Lubelski, \textit{Prace Matematyczne -- Fizyczne}, T. 42.}

% === TRANSCRIPTION NOTES ===
% This chunk covers p100.png through p118.png (original pages "91-א" and 92-109).
% Pages 100-103 (original pages 91-א, 92, 93 partial) contain heavily degraded
%   introductory prose (Chapter 7 opening, section 1 "Introduction") where the
%   typewriter scan quality made several words genuinely ambiguous. These
%   passages are marked inline with "% UNCERTAIN wording" comments. The overall
%   meaning was reconstructed from context (generalizing Pell's equation with
%   respect to exponents, referencing Dirichlet's unit theory, etc.) but exact
%   word choices in a few sentences could not be confirmed with confidence.
% Individual words flagged "% ILLEGIBLE:" appear on original pages 95, 96 (twice),
%   98 (three times), 99, 102, 103, 104, 105, 106 -- these are short connector
%   words or a technical term/reference name (page 104, "פוארקדה") that could not
%   be read with certainty from the scan; the surrounding mathematical content
%   (equations, formulas, theorem statements) is transcribed with high confidence
%   since it is much more legible than the prose.
% Script/curly letters: the cursive symbols used for epsilon/eta/zeta/xi/tau/theta
%   throughout pages 94-106 were rendered as \varepsilon, \eta, \zeta, \xi, \tau,
%   \vartheta per the style guide; the specific glyph shapes in the original are
%   somewhat ambiguous cursive marks and could plausibly be alternate script
%   letters, but the assignments are internally consistent with the algebra.
% The plain typed lines "אוניברסיטה בר אילן / הספריה למתימטיקה" appearing at the
%   bottom of original pages 101 and 109, and the angled ink library stamp on
%   page 109 (p118.png), were treated as library markings per instructions and
%   omitted from the transcription.
% The bibliography heading on p116.png appeared in the scan with unusual letter
%   spacing/possible extra character; per task instructions the standard spelling
%   \section*{ביבליוגרפיה} was used regardless.
% p117.png (original page 108) has no clearly visible page-number footer in the
%   scan; its page number (108) was inferred from sequential position between
%   p116 (107) and p118 (109).
% Several journal volume/page numbers in the bibliography (pages 107-109) are
%   faint or partly cut off in the scan (e.g. "Vierteljahrschrift ... Zürich, 32,
%   1887", "Archiv Math. Phys., 52, 1871", author name "G. Fontané" partially
%   legible as "Fjton..."), and author "O. Perron" (partially legible as "O. Per...")
%   were reconstructed from standard bibliographic knowledge of this literature
%   combined with the legible fragments; flagged here for verification.
